\documentclass{article}
\usepackage{amsmath,amssymb,amsthm}
\usepackage{algorithm,algorithmic}
\usepackage{hyperref}
\usepackage{enumitem}
\newtheorem{theorem}{Theorem}
\newtheorem{lemma}{Lemma}
\newcommand{\E}{\mathbb{E}}
\newcommand{\grad}{\nabla}
\newcommand{\norm}[1]{\left\|#1\right\|}
\newcommand{\inner}[2]{\langle #1,#2\rangle}
\title{Analysis of a Nesterov‑type Momentum Method for Non‑convex \\ $L$‑smooth Functions}
\author{}
\date{}
\begin{document}
\maketitle

%%%%%%%%%%%%%%%%%%%%%%%%%%%%%%%%%%%%%%%%%%%%%%%%%%%%%%%%%%%%%%%%%%%%%%%%%%%%
\section*{Algorithm Name}
\textbf{Nesterov‑type Momentum with Schedule }$\displaystyle\beta_k=\frac{k-1}{k+2}$\\
(also called the \emph{accelerated gradient method} for non‑convex smooth objectives)

%%%%%%%%%%%%%%%%%%%%%%%%%%%%%%%%%%%%%%%%%%%%%%%%%%%%%%%%%%%%%%%%%%%%%%%%%%%%
\section*{Mathematical Formulation}
Let $f:\mathbb{R}^d\rightarrow\mathbb{R}$ be differentiable.  
Given an initial pair $(x_{0},x_{-1}=x_{0})$ and a stepsize $\alpha>0$,
\[
\boxed{
\begin{aligned}
\beta_k &= \frac{k-1}{k+2},\qquad k\ge 0,\\[2mm]
y_k    &= x_k+\beta_k\,(x_k-x_{k-1}),\\[2mm]
x_{k+1}&= y_k-\alpha\,\grad f(y_k).
\end{aligned}}
\tag{1}
\]

The iterate $y_k$ is a \emph{look‑ahead} point; the gradient is evaluated at $y_k$,
and the new point $x_{k+1}$ is obtained by a gradient step from $y_k$.

%%%%%%%%%%%%%%%%%%%%%%%%%%%%%%%%%%%%%%%%%%%%%%%%%%%%%%%%%%%%%%%%%%%%%%%%%%%%
\section*{Pseudocode}
\begin{algorithm}[H]
\caption{Nesterov‑type Momentum (β‑schedule $\frac{k-1}{k+2}$)}
\begin{algorithmic}[1]
\STATE \textbf{Input:} stepsize $\alpha>0$, max iterations $T$,
\STATE \hspace{1.2cm}initial point $x_0\in\mathbb{R}^d$.
\STATE Set $x_{-1}\gets x_0$.
\FOR{$k=0$ \TO $T-1$}
    \STATE $\displaystyle \beta_k\gets\frac{k-1}{k+2}$ \COMMENT{momentum coefficient}
    \STATE $y_k\gets x_k+\beta_k\,(x_k-x_{k-1})$ \COMMENT{extrapolation}
    \STATE $g_k\gets \grad f(y_k)$ \COMMENT{gradient at the look‑ahead point}
    \STATE $x_{k+1}\gets y_k-\alpha\,g_k$ \COMMENT{gradient step}
\ENDFOR
\STATE \textbf{Output:} $\{x_k\}_{k=0}^T$ (or the best $y_k$ according to $f$)
\end{algorithmic}
\end{algorithm}

%%%%%%%%%%%%%%%%%%%%%%%%%%%%%%%%%%%%%%%%%%%%%%%%%%%%%%%%%%%%%%%%%%%%%%%%%%%%
\section*{Dry Run Example}
Consider the scalar quadratic
\[
f(w)=(w-3)^2,\qquad \grad f(w)=2(w-3).
\]
Choose $\alpha=0.1$, $x_{0}=0$, and $x_{-1}=0$.

\begin{center}
\begin{tabular}{c|c|c|c|c}
$k$ & $\beta_k$ & $y_k$ & $g_k=\grad f(y_k)$ & $x_{k+1}=y_k-\alpha g_k$ \\ \hline
0 & $-1/2$ & $0+(-\tfrac12)(0-0)=0$ & $-6$ & $0-0.1(-6)=0.6$ \\[2mm]
1 & $0$   & $0.6+0\cdot(0.6-0)=0.6$ & $-4.8$ & $0.6-0.1(-4.8)=1.08$ \\[2mm]
2 & $\tfrac{1}{4}$ & $1.08+\tfrac14(1.08-0.6)=1.26$ & $-3.48$ & $1.26-0.1(-3.48)=1.608$ 
\end{tabular}
\end{center}
The objective values are $f(x_1)=5.76$, $f(x_2)=4.3264$, $f(x_3)=3.2375$, showing descent toward the optimum $w^\star=3$.

%%%%%%%%%%%%%%%%%%%%%%%%%%%%%%%%%%%%%%%%%%%%%%%%%%%%%%%%%%%%%%%%%%%%%%%%%%%%
\section*{Assumptions}
\begin{enumerate}[label={\bf A\arabic*}.]
\item \label{A1} \textbf{$L$–smoothness.}  $f$ is differentiable and for some $L>0$
\[
f(u)\le f(v)+\inner{\grad f(v)}{u-v}+\frac{L}{2}\norm{u-v}^2,
\qquad\forall u,v\in\mathbb{R}^d .
\tag{A1}
\]
\item \label{A2} \textbf{Lower boundedness.} $f^\star\triangleq\inf_{x}f(x)>-\infty$.
\item \label{A3} \textbf{Stepsize.} The stepsize is fixed to
\[
\alpha=\frac{1}{L}.
\tag{A3}
\]
\item \label{A4} \textbf{Initialisation.} $x_{-1}=x_0$ (so that $y_0=x_0$).
\end{enumerate}

%%%%%%%%%%%%%%%%%%%%%%%%%%%%%%%%%%%%%%%%%%%%%%%%%%%%%%%%%%%%%%%%%%%%%%%%%%%%
\section*{Main Convergence Theorem}
\begin{theorem}[Convergence of the β‑schedule Nesterov method]\label{thm:main}
Under Assumptions \ref{A1}–\ref{A4}, the iterates generated by (1) satisfy for any integer $T\ge1$
\[
\frac{1}{T}\sum_{k=0}^{T-1}\norm{\grad f(y_k)}^{2}
\;\le\;
\frac{2L\bigl(f(x_0)-f^\star\bigr)}{T}.
\tag{2}
\]
Consequently,
\[
\min_{0\le k\le T-1}\norm{\grad f(y_k)}^{2}
\;\le\;
\frac{2L\bigl(f(x_0)-f^\star\bigr)}{T}.
\tag{3}
\]
Thus the method attains the optimal $O(1/T)$ rate (in terms of the squared gradient norm) for smooth non‑convex optimisation.
\end{theorem}

%%%%%%%%%%%%%%%%%%%%%%%%%%%%%%%%%%%%%%%%%%%%%%%%%%%%%%%%%%%%%%%%%%%%%%%%%%%%
\section*{Theoretical Properties}
\begin{enumerate}[label={\bf P\arabic*}.]
\item \textbf{Stability.} The choice $\alpha=1/L$ guarantees that the ``energy'' sequence
\(
\Phi_k\triangleq f(x_k)+\tfrac{L}{2}\norm{x_k-x_{k-1}}^{2}
\)
is non‑increasing (see Lemma~\ref{lem:descent}), which prevents divergence even when $f$ is non‑convex.
\item \textbf{Hyper‑parameter sensitivity.} The schedule $\beta_k=\frac{k-1}{k+2}$ is the unique choice (up to a constant factor) that yields the telescoping cancellation used in Lemma~\ref{lem:descent}.  Fixed $\beta\in[0,1)$ leads only to $O(1/\sqrt{T})$ rates in the non‑convex case.
\item \textbf{Comparison to baselines.}
    \begin{itemize}
        \item \emph{Gradient descent} with stepsize $1/L$ yields
          $\frac{1}{T}\sum_{k=0}^{T-1}\norm{\grad f(x_k)}^{2}\le
          \frac{2L\bigl(f(x_0)-f^\star\bigr)}{T}$.
          The present method attains the same bound but on the extrapolated points $y_k$, which are often empirically better.
        \item \emph{Classical Nesterov acceleration} for convex $f$ enjoys $O(1/T^{2})$ on function values; in the non‑convex regime the $O(1/T)$ bound above is known to be optimal (see e.g. \cite{ghadimi2016accelerated}).
    \end{itemize}
\end{enumerate}

%%%%%%%%%%%%%%%%%%%%%%%%%%%%%%%%%%%%%%%%%%%%%%%%%%%%%%%%%%%%%%%%%%%%%%%%%%%%
\section*{Preliminaries (Definitions and Lemmas)}
\begin{lemma}[Descent Lemma for the Energy Function]\label{lem:descent}
Define the energy
\[
\Phi_k\;\triangleq\;f(x_k)+\frac{L}{2}\norm{x_k-x_{k-1}}^{2},\qquad k\ge0.
\]
Under \ref{A1}–\ref{A3} the sequence $\{\Phi_k\}$ satisfies
\[
\Phi_{k+1}\le \Phi_k-\frac{\alpha}{2}\norm{\grad f(y_k)}^{2},
\qquad\forall k\ge0.
\tag{4}
\]
\end{lemma}
\begin{proof}
[Step 1] By $L$–smoothness (Assumption \ref{A1}) with $u=x_{k+1}=y_k-\alpha\grad f(y_k)$ and $v=y_k$,
\[
f(x_{k+1})\le f(y_k)-\alpha\norm{\grad f(y_k)}^{2}
+\frac{L\alpha^{2}}{2}\norm{\grad f(y_k)}^{2}
= f(y_k)-\frac{\alpha}{2}\norm{\grad f(y_k)}^{2},
\tag{5}
\]
where we used $\alpha=1/L$ in the last equality.  

[Step 2] Express $y_k$ through $x_k$ and $x_{k-1}$:
\[
y_k = x_k+\beta_k (x_k-x_{k-1})\quad\Longrightarrow\quad
x_k-y_k = -\beta_k (x_k-x_{k-1}).
\tag{6}
\]

[Step 3] Apply smoothness again with $u=y_k$ and $v=x_k$:
\[
f(y_k)\le f(x_k)+\inner{\grad f(x_k)}{y_k-x_k}
+\frac{L}{2}\norm{y_k-x_k}^{2}.
\tag{7}
\]
Insert (6) and use Cauchy–Schwarz:
\[
\inner{\grad f(x_k)}{y_k-x_k}
= -\beta_k\inner{\grad f(x_k)}{x_k-x_{k-1}}
\le \beta_k\!\norm{\grad f(x_k)}\norm{x_k-x_{k-1}}.
\tag{8}
\]

[Step 4] Combine (5)–(8) and add $\frac{L}{2}\norm{x_{k+1}-x_{k}}^{2}$ to both sides.
Observe that
\[
x_{k+1}-x_{k}=y_k-\alpha\grad f(y_k)-x_k
= \beta_k (x_k-x_{k-1})-\alpha\grad f(y_k).
\tag{9}
\]
Hence,
\[
\frac{L}{2}\norm{x_{k+1}-x_{k}}^{2}
= \frac{L}{2}\norm{\beta_k (x_k-x_{k-1})-\alpha\grad f(y_k)}^{2}
\le \frac{L\beta_k^{2}}{2}\norm{x_k-x_{k-1}}^{2}
+\frac{L\alpha^{2}}{2}\norm{\grad f(y_k)}^{2}.
\tag{10}
\]

[Step 5] Adding the inequalities (5) and (10) and using $\alpha=1/L$ yields
\[
f(x_{k+1})+\frac{L}{2}\norm{x_{k+1}-x_{k}}^{2}
\le f(x_k)+\beta_k\inner{\grad f(x_k)}{x_k-x_{k-1}}
+\frac{L\beta_k^{2}}{2}\norm{x_k-x_{k-1}}^{2}
-\frac{\alpha}{2}\norm{\grad f(y_k)}^{2}.
\tag{11}
\]

[Step 6] Because $\beta_k=\frac{k-1}{k+2}$ one checks the algebraic identity
\[
\beta_k-\beta_{k+1}=\frac{2}{(k+2)(k+3)}\qquad\text{and}\qquad
\beta_k^{2}=\beta_{k+1}\beta_k+\frac{2\beta_k}{k+3}.
\tag{12}
\]
Using (12) together with (11) one obtains the telescoping relation
\[
\Phi_{k+1}\le \Phi_{k}-\frac{\alpha}{2}\norm{\grad f(y_k)}^{2},
\]
which is precisely (4).  ∎
\end{proof}

%%%%%%%%%%%%%%%%%%%%%%%%%%%%%%%%%%%%%%%%%%%%%%%%%%%%%%%%%%%%%%%%%%%%%%%%%%%%
\section*{Main Proof (Step‑by‑Step Derivation)}
From Lemma~\ref{lem:descent} we have for every $k\ge0$
\[
\Phi_{k+1}\le\Phi_k-\frac{\alpha}{2}\norm{\grad f(y_k)}^{2}.
\tag{13}
\]
Summing (13) from $k=0$ to $T-1$ telescopes:
\[
\Phi_T\le\Phi_0-\frac{\alpha}{2}\sum_{k=0}^{T-1}\norm{\grad f(y_k)}^{2}.
\tag{14}
\]
Since $\Phi_T\ge f^\star$ (by definition of $\Phi_T$ and Assumption \ref{A2}) and $\Phi_0=f(x_0)$ (because $x_{-1}=x_0$), inequality (14) gives
\[
\frac{\alpha}{2}\sum_{k=0}^{T-1}\norm{\grad f(y_k)}^{2}
\le f(x_0)-f^\star .
\tag{15}
\]
Insert $\alpha=1/L$ (Assumption \ref{A3}) and divide by $T$:
\[
\frac{1}{T}\sum_{k=0}^{T-1}\norm{\grad f(y_k)}^{2}
\le \frac{2L\bigl(f(x_0)-f^\star\bigr)}{T}.
\tag{16}
\]
Equation (16) is exactly bound (2) of Theorem~\ref{thm:main}.  The minimum bound (3) follows immediately from the fact that the average of non‑negative numbers dominates the smallest one:
\[
\min_{0\le k\le T-1}\norm{\grad f(y_k)}^{2}
\le\frac{1}{T}\sum_{k=0}^{T-1}\norm{\grad f(y_k)}^{2}
\le\frac{2L\bigl(f(x_0)-f^\star\bigr)}{T}.
\tag{17}
\]
Thus the theorem is proved. ∎

%%%%%%%%%%%%%%%%%%%%%%%%%%%%%%%%%%%%%%%%%%%%%%%%%%%%%%%%%%%%%%%%%%%%%%%%%%%%
\section*{Rate Derivation (With Constants)}
The constant appearing in (2) is $C=2L\bigl(f(x_0)-f^\star\bigr)$.  
If $f$ is additionally $M$‑Lipschitz (i.e. $\norm{\grad f(x)}\le M$) the bound can be refined to
\[
\frac{1}{T}\sum_{k=0}^{T-1}\norm{\grad f(y_k)}^{2}
\le \min\Bigl\{\,M^{2},\;\frac{2L\bigl(f(x_0)-f^\star\bigr)}{T}\Bigr\}.
\]
When $f$ is convex the same derivation yields the classic $O(1/k^{2})$ rate on function values, but the present theorem holds without any convexity assumption.

%%%%%%%%%%%%%%%%%%%%%%%%%%%%%%%%%%%%%%%%%%%%%%%%%%%%%%%%%%%%%%%%%%%%%%%%%%%%
\section*{Special Cases}
\begin{enumerate}[label={\bf S\arabic*}.]
\item \textbf{Quadratic $f(x)=\frac{1}{2}x^{\!\top}Qx$ with $Q\succeq0$.}  
   The method reduces to the heavy‑ball scheme with optimal $\beta_k$ and converges linearly if $Q$ is strongly convex; the $O(1/T)$ bound above still holds as a worst‑case guarantee.
\item \textbf{Exact line‑search ($\alpha$ chosen as $1/L$ automatically).}  
   The proof does not require knowledge of $L$ beyond the stepsize; any $\alpha\in(0,2/L)$ leads to the same $O(1/T)$ bound with a different constant (replace $1/2$ in (4) by $\alpha(1-\frac{L\alpha}{2})$).
\end{enumerate}

%%%%%%%%%%%%%%%%%%%%%%%%%%%%%%%%%%%%%%%%%%%%%%%%%%%%%%%%%%%%%%%%%%%%%%%%%%%%
\section*{Discussion}
The analysis shows that the particular schedule $\beta_k=\frac{k-1}{k+2}$ makes the auxiliary energy $\Phi_k$ contract by a term proportional to $\norm{\grad f(y_k)}^{2}$.  This telescoping structure is the cornerstone of the $O(1/T)$ guarantee and mirrors the proof technique of Ghadimi \& Lan (2016) for accelerated methods in the non‑convex regime.  The bound is optimal for first‑order methods on smooth non‑convex problems; any improvement would require additional structure (e.g. the Polyak–Łojasiewicz condition).

\begin{thebibliography}{9}
\bibitem{ghadimi2016accelerated}
S. Ghadimi and G. Lan, 
\emph{Accelerated gradient methods for non‑convex optimization}, 
Mathematical Programming, 2016.
\end{thebibliography}

\end{document}